\section{Conclusion}\label{sec:conclusion}

This paper assembles a single quantitative-finance artifact around the development trajectory of 23 emerging economies. The work spans three deliberately distinct lenses: causal macro identification, sovereign credit pricing and emerging-market factor investing. The composite scoring machine sits underneath all three; the data layer underneath the scoring; and the panel framework underneath the data. Every figure and table in this document is generated from a single Python pipeline that re-runs in roughly two minutes on a warm cache. The composite ranks Singapore first and Niger last, which is the right answer for a development-KPI score and is the answer that survives the per-capita and logarithmic transforms that the size of China and India would otherwise dominate.

The substantive contributions are three.

\emph{First}, the causal block delivers more than the LP IRFs alone might suggest. The Hamilton-cyclical innovations of slow-moving development indicators do not drive LP point estimates outside conventional confidence bands. But the same panel under the Dumitrescu-Hurlin panel-Granger test --- which averages country-level Wald statistics rather than imposing a homogeneous $\beta_h$ --- finds that energy shocks Granger-cause poverty reduction and health shocks Granger-cause Gini reduction at conventional significance ($p \le 0.05$), while no sector Granger-causes log GDP per capita. The take-away is that KPI improvements feed the distributional tail (poverty headcount, Gini) before they show up in aggregate GDP --- a substantive finding consistent with the development-economics intuition that infrastructure, sanitation and basic-health gains raise the floor rather than the mean.

\emph{Second}, the credit block is the empirical headline. A pooled ordered probit of the year-end S\&P rating on sector scores alone delivers an out-of-sample MAE of 1.84 notches and an AUC of 0.81 for the investment-grade binary, against the 0.73 AUC of the standard macro-only baseline. A Bayesian hierarchical ordered-logistic counterpart converges cleanly (NUTS, 4 chains, 1\,500 draws each, $\sim$100 seconds) and localises the signal to two sectors with credibly non-zero population loadings: energy ($\mu_\beta = +0.92$, 90~\% HDI $[+0.10,\,+1.74]$) and health ($\mu_\beta = -0.92$, HDI $[-1.38,\,-0.45]$); the four other sectors have HDIs that cover zero. The implication for a fixed-income reader is direct: a development-KPI overlay can improve a rating-driven mandate's discrimination without adding any of the conventional macro inputs, and the parsimonious version that loads on energy and health alone is supported by the Bayesian posterior.

\emph{Third}, the portfolio block delivers \emph{negative} risk-adjusted alpha relative to a random-portfolio benchmark. The KPI-sorted top-quintile EM portfolio realises a Sharpe of 0.15 against the MSCI EM benchmark's 0.12 over 2007--2024 --- a marginal lift in headline terms --- but the same Sharpe sits at the $0^{\text{th}}$ percentile of the 5\,000-random-basket distribution (mean random Sharpe $0.28$). The KPI tilt mis-allocates by roughly $0.13$ Sharpe units relative to an equal-weighted random EM basket. The long-short variant is a disaster (Sharpe $-0.41$, max DD 94~\%). The result is consistent with the slow-moving nature of the KPI signal: information about fundamentals shows up where the underlying instrument prices slowly (sovereign ratings), not where it prices fast (monthly equity).

The three results are not contradictions; they are complements. KPIs are slow-moving, so they price into slow-moving instruments. The natural follow-up extends the portfolio exercise from country ETFs to EMBI sovereign credit spreads, where the credit-pricing finding of Section~\ref{sec:partII} would translate more directly. That extension is the next milestone of this work.

\subsection*{Limitations}

The 30-year hand-compiled S\&P rating CSV deserves an external audit before any operational use --- the source-of-record is the S\&P sovereign archive and individual press releases, but the year-end snapshots have not been formally back-checked here. PISA scores are sparse for eight of the focal 13 countries and were demoted to a supplementary status. The IEA-balance and UN-Comtrade ingestion paths are stubbed; these gates are documented in \texttt{src/newperformers/io/}. None of these limitations changes the qualitative reading of the three blocks above.
